% Clase 28 — la inconsistencia de Ampère: un cable con corriente I(t) llega a
% un condensador. Con la misma curva C, la superficie S1 la cruza el cable
% (da I) y la S2 pasa entre las placas (daría 0). La diferencia es la carga que
% se acumula en la placa.
\begin{tikzpicture}[scale=1]

  % cable y placas
  \draw[figblue, line width=1.6pt] (-3.0,0) -- (1.2,0);
  \draw[figblue, line width=1.6pt] (2.4,0) -- (4.2,0);
  \draw[figblue, line width=2.2pt] (1.2,-1.3) -- (1.2,1.3);
  \draw[figblue, line width=2.2pt] (2.4,-1.3) -- (2.4,1.3);
  \node[figlblaux, text=figred, anchor=south] at (1.0,1.95) {$+Q(t)$};
  \node[figlblaux, text=figblue, anchor=south] at (2.6,1.95) {$-Q(t)$};
  \draw[figvecres] (-2.6,0.25) -- (-1.8,0.25);
  \node[figlbl, text=figred] at (-2.2,0.6) {$I(t)$};

  % curva C (vista de costado: una elipse alrededor del cable)
  \draw[figred, line width=1pt] (-0.6,0) ellipse (0.18 and 0.95);
  \node[figlbl, text=figred] at (-0.6,-1.25) {$C$};

  % S1: disco plano en la curva
  \fill[figamber!25, opacity=0.8] (-0.6,0) ellipse (0.18 and 0.95);
  \node[figlbl, text=figamber] at (-1.15,0.95) {$S_1$};

  % S2: bolsa que rodea la placa +Q y se cierra entre las placas
  \draw[figamber, line width=0.9pt, dashed]
    (-0.6,0.95) .. controls (0.2,1.75) and (1.8,1.9) .. (1.8,0.9)
               -- (1.8,-0.9)
               .. controls (1.8,-1.9) and (0.2,-1.75) .. (-0.6,-0.95);
  \node[figlbl, text=figamber, fill=white, inner sep=1pt] at (1.8,0.45) {$S_2$};

\end{tikzpicture}
